\documentclass[10pt,a4paper]{amsart}
\usepackage[margin=19mm]{geometry}
\usepackage{amsmath,amssymb,mathtools}
\usepackage[scr=rsfs,cal=euler]{mathalfa}
\usepackage{xfrac}
\usepackage[unicode,hidelinks,bookmarks=false]{hyperref}
\newcommand{\ie}{i.e.\ }
\newcommand{\cf}{cf.\ }
\newcommand{\resp}{respectively\ }
\newcommand{\Subsection}{\subsection}
\providecommand{\nobreakdash}{\nobreak\hbox{-}}
\setlength{\emergencystretch}{2em}
\multlinegap0pt
\usepackage{mathrsfs}
\usepackage{fullpage}
\usepackage{enumitem}
\newtheorem{theorem}{Theorem}[section]
\newtheorem{proposition}[theorem]{Proposition}
\newtheorem{lemma}[theorem]{Lemma}
\newtheorem{corollary}[theorem]{Corollary}
\newtheorem{definition}[theorem]{Definition}
\newtheorem{example}[theorem]{Example}
\newtheorem{conjecture}[theorem]{Conjecture}
\newtheorem{problem}[theorem]{Problem}
\newtheorem{remark}[theorem]{Remark}
\hypersetup{
	colorlinks=true,
	linkcolor=blue,
	citecolor=blue,
	urlcolor=blue
}
\begin{document}
\title[Cohen-Macaulay higher conormal and K"ahler differential modu]{Cohen-Macaulay higher conormal and K\"ahler differential modules of squarefree monomial ideals}
\author{T\`ai Huy H\`a; Nguyen Cong Minh}
\date{}
\maketitle
\noindent\emph{Source and licence.} Cohen-Macaulay higher conormal and K\"ahler differential modules of squarefree monomial ideals. Version arXiv:2609.18071v1 (CC BY 4.0). This material is distributed under the Creative Commons Attribution 4.0 International licence, http://creativecommons.org/licenses/by/4.0/, as recorded in the arXiv OAI licence metadata for this exact version and shown on the abstract page https://arxiv.org/abs/2609.18071v1. Full rights notice: licenses/arxiv-2609.18071-CC-BY-4.0.txt.\par\smallskip
\noindent\textit{Editorial scope.} Counted unit: the original \texttt{theorem} environment \texttt{CM-edge-ideals} at source lines 584--595 together with its complete proof at lines 597--706; the delivered document keeps source lines 340--707 so that every referenced label and every citation resolves inside it. Native editable source is the paper's own concatenated TeX; the AMS wrapper is editorial formatting only. No publisher class, upstream script or unrelated artwork is redistributed.
\begin{lemma}\label{lem:negative-link}
For every monomial ideal $J\subseteq S$ and every $\mathbf a\in\mathbb Z^n$,
$$
\Delta_{\mathbf a}(J)
=
\operatorname{link}_{\Delta_{\mathbf a^+}(J)}(G_{\mathbf a}).
$$
\end{lemma}

\begin{proof}
A face $F\subseteq[n]\setminus G_{\mathbf a}$ belongs to the left-hand side
if and only if, after inverting the variables in $G_{\mathbf a}$, the monomial
$x^{\mathbf a^+}$ does not belong to $J S_{x_{F\cup G_{\mathbf a}}}$. This is
exactly the condition that $F\cup G_{\mathbf a}$ belongs to
$\Delta_{\mathbf a^+}(J)$.
\end{proof}

If $J\subseteq I$ are monomial ideals, then
$\Delta_{\mathbf a}(I)\subseteq\Delta_{\mathbf a}(J)$, and hence
$(\Delta_{\mathbf a}(J),\Delta_{\mathbf a}(I))$ is a relative simplicial
complex. We use the following relative Hochster--Takayama formula from
\cite[Theorem~3.2]{HaMinh}: whenever
$G_{\mathbf a}\in\Delta(\sqrt J)$,
$$
H^i_{\mathfrak m}(I/J)_{\mathbf a}
\cong
\widetilde H^{i-|G_{\mathbf a}|-1}
\bigl(\Delta_{\mathbf a}(J),\Delta_{\mathbf a}(I);k\bigr).
$$
Over the field $k$, nonvanishing of reduced relative homology and cohomology in
degree zero are equivalent. We shall therefore use $\widetilde H_0$ in the
combinatorial arguments below. We will also use the following elementary
relative-homology observation.

\begin{lemma}\label{lem:relative-zero}
	Let $X=A\cup B$ be a simplicial complex such that $A$ and $B$ each contain a
	vertex and $A\cap B=\{\emptyset\}$. If $C\subseteq A$ is a simplicial
	subcomplex, then $\widetilde H_0(X,C;k)\neq0$.
\end{lemma}

\begin{proof} This observation is taken from \cite[Lemma 3.1]{MinhNakamura}. For the reader's convenience, we provide a direct proof here. No nonempty face of $X$ meets both $A$ and $B$, so no connected component of
	$X$ meets both subcomplexes. If $C$ has no vertex, then $X$ is disconnected and
	$\widetilde H_0(X;k)\neq0$. Since $\widetilde H_0(C;k)=0$, the long exact
	sequence of the pair gives an injection
	$\widetilde H_0(X;k)\to\widetilde H_0(X,C;k)$, so the latter group is nonzero.
	If $C$ has a vertex, any connected component contained in $B$ is disjoint from
	$C$ and determines a nonzero class in $H_0(X,C;k)$. Since $C$ is nonempty in
	this case, ordinary and reduced relative homology agree in degree zero.
\end{proof}

\section{Ordinary higher conormal filtration}\label{sec:ordinary}

For the rest of the paper, unless stated otherwise, let
$I=I_\Delta\subsetneq S$ be a nonzero squarefree monomial ideal. We first
record the support calculation for arbitrary windows in both filtrations.

\begin{lemma}\label{lem:support}
Let $1\le r<q$. Then
$$
\operatorname{Supp}(I^r/I^q)
=
\operatorname{Supp}(I^{(r)}/I^{(q)})
=
V(I).
$$
Consequently,
$\dim(I^r/I^q)=\dim(I^{(r)}/I^{(q)})=\dim(S/I)$.
\end{lemma}

\begin{proof} Let $Q \in \operatorname{Spec } S \setminus V(I)$. Then, $I \not\subseteq Q$. Choose $f\in I\setminus Q$. Since
$f^r\in I^r\subseteq I^{(r)}$ and $f^q\in I^q\subseteq I^{(q)}$,
localization at $Q$ makes each of $I^r$, $I^q$, $I^{(r)}$, and $I^{(q)}$
equal to $S_Q$. Thus, both quotients vanish after localization at $Q$.

Conversely, let $P\in\operatorname{Min}(I)$. Since $I$ is radical,
$I_P=PS_P$, while symbolic powers commute with localization at a minimal
prime. Hence,
$$
(I^r/I^q)_P\cong P^rS_P/P^qS_P
$$
and
$$
(I^{(r)}/I^{(q)})_P\cong P^rS_P/P^qS_P.
$$
These modules are nonzero because $PS_P$ is a nonzero maximal ideal of the regular local ring $S_P$ and its powers are strict. Thus, every minimal prime of $I$
belongs to both supports. Since support is specialization closed, both supports
are $V(I)$.
\end{proof}

The only pair $1\le r<q$ with $q<3$ is $(r,q)=(1,2)$. At this exceptional
level, ordinary Cohen--Macaulayness forces the ordinary and symbolic squares
to coincide.

\begin{proposition}\label{prop:second-symbolic}
Let $I\subsetneq S$ be a nonzero squarefree monomial ideal. If $I/I^2$ is
Cohen--Macaulay, then
$$
I^2=I^{(2)}.
$$
In particular, $I/I^{(2)}$ is Cohen--Macaulay.
\end{proposition}

\begin{proof}
Set $M=I/I^2$ and $d=\dim(S/I)$. By Lemma~\ref{lem:support},
$\operatorname{Supp}(M)=V(I)$ and $\dim M=d$. Since $M$ is a finitely
generated Cohen--Macaulay $S$-module, it is unmixed; equivalently,
$\operatorname{Ass}_S(M)=\operatorname{Assh}_S(M)$. Hence,
$$
\dim S/Q=d\text{ for every }Q\in\operatorname{Ass}_S(M).
$$
Every minimal prime of $I$ is a minimal prime of $\operatorname{Supp}(M)$,
and hence belongs to $\operatorname{Ass}_S(M)$. We claim that
$$
\operatorname{Ass}_S(M)=\operatorname{Min}(I).
$$
Indeed, let $Q\in\operatorname{Ass}_S(M)$. Since $Q\in V(I)$, there is
$P\in\operatorname{Min}(I)$ with $P\subseteq Q$. Both $P$ and $Q$ belong
to $\operatorname{Ass}_S(M)$, so $\dim S/P=\dim S/Q=d$. A strict
inclusion of primes in $S$ strictly decreases dimension, and therefore $P=Q$.

Set $T=I^{(2)}/I^2\subseteq M$. For every $P\in\operatorname{Min}(I)$,
radicality gives $I_P=PS_P$, and symbolic localization gives
$$
(I^{(2)})_P=P^2S_P=(I_P)^2=(I^2)_P.
$$
Thus, $T_P=0$ for every $P\in\operatorname{Min}(I)$. If $T\neq0$, choose
$Q\in\operatorname{Ass}_S(T)$. Since $T\subseteq M$,
$\operatorname{Ass}_S(T)\subseteq\operatorname{Ass}_S(M)=\operatorname{Min}(I)$.
Hence, $Q$ is a minimal prime of $I$, but $Q\in\operatorname{Ass}_S(T)$
implies $T_Q\neq0$, a contradiction. Therefore $T=0$, so
$I^2=I^{(2)}$. The final assertion follows because then
$I/I^{(2)}=I/I^2$.
\end{proof}

For upper exponent at least three, the following local configuration is the
basic obstruction for every ordinary window.

\begin{lemma}\label{lem:three-vertex-obstruction}
Let $1\le r<q$ with $q\ge3$. Suppose that, after relabeling three vertices
of $\Delta$,
$$
\{1\},\{2,3\}\in\Delta
\text{ and }
\{1,2\},\{1,3\}\notin\Delta.
$$
Then, $I^r/I^q$ is not Cohen--Macaulay.
\end{lemma}

\begin{proof}
Set $\mathbf b=(q-1)\mathbf e_1+\mathbf e_2+\mathbf e_3$. We first claim
that
$$
\Delta_{\mathbf b}(I^q)
=
\operatorname{star}_\Delta(1)
\cup
\operatorname{star}_\Delta(\{2,3\}).
$$
Let $F\in\operatorname{star}_\Delta(1)$. If
$x^{\mathbf b}=x_1^{q-1}x_2x_3$ belonged to $I^qS_{x_F}$, then there would
be $q$ minimal monomial generators of $I$ whose product divides
$x^{\mathbf b}$ after the variables indexed by $F$ are declared units.
Since $F\in\Delta$, none of these generators can become a unit. Moreover,
the noninvertible part of each generator must be divisible by $x_2$ or $x_3$;
otherwise its support would give a nonface contained in $F\cup\{1\}$. This is impossible because
$x^{\mathbf b}$ has total degree two in $x_2,x_3$ and $q\ge3$. Hence
$F\in\Delta_{\mathbf b}(I^q)$.

If $F\in\operatorname{star}_\Delta(\{2,3\})$, the same argument shows
that the noninvertible part of each of the $q$ factors must be divisible by
$x_1$. Their product would then be divisible by $x_1^q$, contradicting the
exponent $q-1$ of $x_1$ in $x^{\mathbf b}$. Thus, the second star is also
contained in $\Delta_{\mathbf b}(I^q)$.

Conversely, suppose that $F$ belongs to neither star. If $F\notin\Delta$,
then $IS_{x_F}=S_{x_F}$, so $F\notin\Delta_{\mathbf b}(I^q)$. Assume
$F\in\Delta$. Since $F\cup\{1\}\notin\Delta$, a minimal nonface
contained in $F\cup\{1\}$ contains $1$, and localization at $F$ gives
$x_1\in IS_{x_F}$. Similarly, since
$F\cup\{2,3\}\notin\Delta$, one of $x_2,x_3,x_2x_3$ belongs to
$IS_{x_F}$. Taking $q-1$ copies of $x_1$ and one copy of this last monomial
shows that $x^{\mathbf b}\in I^qS_{x_F}$. This proves the claimed equality.

We next claim that
$$
\Delta_{\mathbf b}(I^r)\subseteq\operatorname{star}_\Delta(1).
$$
Indeed, let $F\notin\operatorname{star}_\Delta(1)$. If
$F\notin\Delta$, then $I^rS_{x_F}=S_{x_F}$. If $F\in\Delta$, the same
minimal-nonface argument gives $x_1\in IS_{x_F}$, so
$x_1^r\in I^rS_{x_F}$. Since $r\le q-1$, the monomial $x_1^r$ divides
$x^{\mathbf b}$. In either case $F\notin\Delta_{\mathbf b}(I^r)$,
proving the inclusion.

Set
$$
L=\operatorname{link}_\Delta(1)
\cap
\operatorname{link}_\Delta(\{2,3\}),
$$
choose a maximal face $U$ of $L$, and put
$\mathbf a=\mathbf b-\sum_{i\in U}\mathbf e_i$. Then
$G_{\mathbf a}=U$ and $\mathbf a^+=\mathbf b$. By
Lemma~\ref{lem:negative-link},
$$
\Delta_{\mathbf a}(I^q)=A\cup B,
$$
where
$$
A=\operatorname{link}_{\operatorname{star}_\Delta(1)}(U)
\text{ and }
B=\operatorname{link}_{\operatorname{star}_\Delta(\{2,3\})}(U).
$$
The complex $A$ contains the vertex $1$, while $B$ contains the edge
$\{2,3\}$. Moreover, $A\cap B=\{\emptyset\}$. Indeed, if $W$ were a
nonempty face of $A\cap B$, then
$U\cup W\cup\{1\}\in\Delta$ forces $2,3\notin W$, while
$U\cup W\cup\{2,3\}\in\Delta$ forces $1\notin W$. Hence,
$U\cup W\in L$, contradicting the maximality of $U$.

The inclusion above and Lemma~\ref{lem:negative-link} give
$\Delta_{\mathbf a}(I^r)\subseteq A$. Therefore,
Lemma~\ref{lem:relative-zero} yields
$$
\widetilde H_0
\bigl(\Delta_{\mathbf a}(I^q),\Delta_{\mathbf a}(I^r);k\bigr)
\neq0.
$$
Since $U\in\Delta=\Delta(\sqrt{I^q})$, the relative
Hochster--Takayama formula applies and gives
$$
H_{\mathfrak m}^{|U|+1}(I^r/I^q)_{\mathbf a}\neq0.
$$
Since $U\cup\{2,3\}\in\Delta$, one has
$|U|+2\le\dim\Delta+1$. By Lemma~\ref{lem:support},
$$
|U|+1\le\dim\Delta<\dim\Delta+1=\dim(I^r/I^q).
$$
Thus, $I^r/I^q$ is not Cohen--Macaulay.
\end{proof}

For edge ideals, the obstruction gives a concrete classification for all
power-filtration windows.


\begin{theorem}\label{CM-edge-ideals}
Let $G$ be a finite simple graph with at least one edge, let
$S=k[x_1,\ldots,x_n]$, and let $I=I(G)$. Fix integers $1\le r<q$ with
$q\ge3$. Then, the following are equivalent:
\begin{enumerate}
\item $I^r/I^q$ is Cohen--Macaulay;
\item $I$ is a complete intersection;
\item after removing isolated vertices, $G$ is a matching.
\end{enumerate}
Consequently, if $I^r/I^q$ is Cohen--Macaulay for one such pair $(r,q)$,
then $I^a/I^b$ is Cohen--Macaulay for every $1\le a<b$.
\end{theorem}

\begin{proof}
Let $\Delta=\operatorname{Ind}(G)$ and set $M=I^r/I^q$. Assume first that
$M$ is Cohen--Macaulay. If some connected component of $G$ were not complete,
it would contain an induced path $2-1-3$. Then,
$$
\{1\},\{2,3\}\in\Delta
\text{ and }
\{1,2\},\{1,3\}\notin\Delta,
$$
contradicting Lemma~\ref{lem:three-vertex-obstruction}. Hence, every connected
component of $G$ is complete. After removing isolated vertices, write
$$
G=K_{r_1}\sqcup\cdots\sqcup K_{r_s},\text{ } r_i\ge2.
$$
All minimal vertex covers of $G$ have cardinality
$\sum_{i=1}^s(r_i-1)$, so $I$ is unmixed. By Lemma~\ref{lem:support},
$\operatorname{Supp}(M)=V(I)$ and $\dim M=\dim(S/I)$. Since $M$ is
Cohen--Macaulay, every associated prime of $M$ has dimension $\dim(S/I)$.
Moreover, every minimal prime of $I$ is a minimal prime of
$\operatorname{Supp}(M)$ and hence belongs to $\operatorname{Ass}_S(M)$.
If $Q\in\operatorname{Ass}_S(M)$ and $P\in\operatorname{Min}(I)$ satisfies
$P\subseteq Q$, then $\dim S/P=\dim S/Q=\dim(S/I)$, so $P=Q$. Therefore,
$$
\operatorname{Ass}_S(M)=\operatorname{Min}(I).
$$

Suppose that some component $C$ is $K_\ell$ with $\ell\ge3$, and choose
distinct vertices $x,y,z$ of $C$. For every other nontrivial complete component, omit
one vertex, and let $P$ be the monomial prime generated by all vertices of
$C$ together with all nonomitted vertices in the other complete components.
No isolated vertex is included in $P$. After localizing at $P$,
$$
I_P=J+L,
$$
where $J=I(K_\ell)$ and $L$ is generated by variables from the other complete
components. The variables supporting $J$ and $L$ are disjoint, and $PS_P$ is
the maximal ideal generated by the variables of $C$ and those generating $L$.

Assume first that $L\neq0$, choose a variable $w\in L$, and set
$$
u=xyz\,w^{q-2}.
$$
Since $xyz\in J$, one has
$u\in JL^{q-2}\subseteq I_P^{q-1}\subseteq I_P^r$. On the other hand,
$u\notin I_P^q$. Indeed,
$$
(J+L)^q=\sum_{a+b=q}J^aL^b.
$$
If $u\in J^aL^b$, then $b\le q-2$, and hence $a\ge2$. Since the
variables supporting $J$ and $L$ are disjoint, this would force the $C$-part
$xyz$ of $u$ to belong to $J^a\subseteq J^2$, which is impossible because
every monomial in $J^2$ has degree at least four.

For every variable $v$ of $C$, one has $vxyz\in J^2$, so
$vu\in J^2L^{q-2}\subseteq I_P^q$. If $t$ is a variable generating $L$,
then $xyz\in J$ and $tw^{q-2}\in L^{q-1}$, so
$tu\in JL^{q-1}\subseteq I_P^q$. Hence,
$I_P^q:u$ contains the maximal ideal $PS_P$. Since $u\notin I_P^q$, this
colon is proper, and therefore
$$
I_P^q:u=PS_P.
$$
Thus, the class of $u$ in $M_P=I_P^r/I_P^q$ has annihilator $PS_P$, so
$PS_P\in\operatorname{Ass}_{S_P}(M_P)$.

If $L=0$, set
$$
u=(xy)^{q-2}xyz.
$$
Since $xyz\in J$, one has $u\in J^{q-1}\subseteq J^r$. Its degree is
$2q-1$, whereas every monomial in $J^q$ has degree at least $2q$, so
$u\notin J^q$. For every variable $v$ of $C$,
$$
xu=(xy)^{q-1}(xz),\text{ } yu=(xy)^{q-1}(yz),
$$
$$
zu=(xy)^{q-2}(xz)(yz),
$$
and, if $v\notin\{x,y,z\}$,
$$
vu=(xy)^{q-2}(xv)(yz).
$$
Thus, $J^q:u=PS_P$, and again $PS_P\in\operatorname{Ass}_{S_P}(M_P)$.

By the localization formula for associated primes, there is
$Q\in\operatorname{Ass}_S(M)$ with $Q\subseteq P$ and
$QS_P=PS_P$. Contracting to $S$ gives $Q=P$, so
$P\in\operatorname{Ass}_S(M)$. But $P$ is not minimal over $I$,
because it contains every vertex of $C$, whereas a minimal vertex cover of
$C=K_\ell$ contains only $\ell-1$ vertices. This contradicts
$\operatorname{Ass}_S(M)=\operatorname{Min}(I)$. Hence, every nontrivial
component of $G$ is $K_2$, proving $(1)\Rightarrow(3)$.

The equivalence of (2) and (3) follows from the characterization of monomial
complete intersections by pairwise disjoint supports of their minimal
generators. If $I$ is a complete intersection, then
$\operatorname{gr}_I(S)\cong\operatorname{Sym}_{S/I}(I/I^2)$; in particular,
each $I^t/I^{t+1}$ is a finite free $S/I$-module. It follows inductively from
$0\to I^t/I^{t+1}\to S/I^{t+1}\to S/I^t\to0$ that $S/I^t$ is
Cohen--Macaulay of dimension $\dim(S/I)$ for every $t\ge1$. The exact
sequence
$$
0\longrightarrow I^r/I^q\longrightarrow S/I^q\longrightarrow S/I^r
\longrightarrow0
$$
together with the Depth Lemma and Lemma~\ref{lem:support} shows that
$I^r/I^q$ is Cohen--Macaulay. This
proves $(2)\Rightarrow(1)$. The final assertion follows because a complete
intersection has Cohen--Macaulay quotients $I^a/I^b$ for all $1\le a<b$.
\end{proof}

\begin{thebibliography}{99}
\bibitem[HaMinh]{HaMinh} T.~H. H\`a and N.~C. Minh, \emph{Relative Hochster--Takayama formula and Cohen--Macaulay monomial ideal quotients}, arXiv:2601.20233 (2026).
\bibitem[MinhNakamura]{MinhNakamura} N. C. Minh and Y. Nakamura, \emph{A note on the k -Buchsbaum property of symbolic powers of Stanley-Reisner ideals}, Tokyo J. Math. \textbf{34} (2011), no. 1, 221–227.
\end{thebibliography}
\end{document}
